\section{Related Work}
\paragraph{Hidden communication between neural agents.}
Abadi and Andersen trained neural networks to protect communications against an adversary with adversarial training~\cite{abadi2016learning}.
Linguistic steganography hides messages in fluent text~\cite{ziegler2019neural,kaptchuk2021meteor}; these methods rely on a shared secret key and a coding scheme,
whereas our key is naturally occurring shared history and nothing is encoded. Motwani et al.\ study secret collusion among AI agents via steganography~\cite{motwani2024secret}.
Roger and Greenblatt study models hiding reasoning in their outputs and defences such as paraphrasing~\cite{roger2023preventing}.
The closest work is \citet{mathew2024hidden}, who train a sender and receiver with reinforcement learning against an overseer that flags suspicious text, with an optional paraphraser.
We differ in three ways: (i) the outsider is a \emph{decoder} built from the same model as the insider listener and told the rules of the game, so its only disadvantage is the missing shared history;
(ii) the key is a freshly generated shared history per episode, written before the secret is drawn; and (iii) the exchange is a two-sided multi-turn dialogue in which the replying insider does not know the secret.
Unlike Mathew et al., our listeners are fixed during training; an outsider that adapts to the speaker is left to future work (Section~\ref{sec:limits}).
\paragraph{Emergent communication.} Lazaridou et al.\ and Foerster et al.\ study how agents develop communication protocols by reinforcement learning~\cite{lazaridou2017multi,foerster2016learning};
we instead start from natural language and ask what can be hidden inside it.
\paragraph{Training methods.} We use GRPO~\cite{shao2024deepseekmath} without a value model and LoRA adapters~\cite{hu2021lora}; the supervised warm start on filtered samples resembles rejection-sampling approaches such as STaR~\cite{zelikman2022star}.
